\documentclass[11pt]{article}

\usepackage[margin=1in]{geometry}
\usepackage{titlesec}
\usepackage{enumitem}
\usepackage{hyperref}
\usepackage{parskip}
\usepackage{amssymb}

\title{Personal Knowledge OS (PKOS) \\ \large Product Requirements Document (PRD)}
\author{}
\date{}

\hypersetup{
    colorlinks=true,
    linkcolor=black,
    urlcolor=blue,
}

\titleformat{\section}{\Large\bfseries}{\thesection}{1em}{}
\titleformat{\subsection}{\large\bfseries}{\thesubsection}{1em}{}
\titleformat{\subsubsection}{\normalsize\bfseries}{\thesubsubsection}{1em}{}

\begin{document}

\maketitle

\section{Vision}

Personal Knowledge OS is an AI-powered learning and knowledge management platform that helps users capture information from multiple sources, organize it into a structured knowledge base, retrieve it intelligently, and retain it through personalized revision recommendations.

The goal is not simply to store information, but to transform scattered content into long-term, reusable knowledge.

\section{Problem Statement}

Modern learners consume information from many sources:

\begin{itemize}[nosep]
    \item Articles
    \item Documentation
    \item YouTube videos
    \item Research papers
    \item Personal notes
    \item Tutorials
\end{itemize}

Over time, several problems emerge:

\begin{enumerate}[nosep]
    \item Information becomes scattered across bookmarks, notes, PDFs, and browser tabs.
    \item Users remember learning something but cannot locate it quickly.
    \item Related concepts remain disconnected.
    \item Valuable knowledge is forgotten due to lack of revision.
    \item Existing note-taking tools focus on storage rather than learning and retention.
\end{enumerate}

As a result, users repeatedly relearn concepts they have already studied.

\section{Target Users}

\subsection{Primary Users}

\subsubsection{Students}

\begin{itemize}[nosep]
    \item Engineering students
    \item Competitive programmers
    \item University learners
\end{itemize}

\textbf{Need:}
\begin{itemize}[nosep]
    \item Organize study materials
    \item Revise concepts efficiently
    \item Track learning progress
\end{itemize}

\subsubsection{Software Engineers}

\textbf{Need:}
\begin{itemize}[nosep]
    \item Store technical knowledge
    \item Quickly retrieve information during development
    \item Build a personal technical knowledge base
\end{itemize}

\subsubsection{AI/ML Learners}

\textbf{Need:}
\begin{itemize}[nosep]
    \item Manage research papers
    \item Connect concepts across topics
    \item Generate summaries and revision material
\end{itemize}

\section{Product Goals}

The platform should help users:

\begin{enumerate}[nosep]
    \item Capture knowledge from any source.
    \item Organize information automatically.
    \item Retrieve knowledge instantly.
    \item Discover relationships between concepts.
    \item Revise information before forgetting it.
    \item Build a long-term personal knowledge repository.
\end{enumerate}

\section{Core Product Pillars}

\subsection{Pillar 1: Knowledge Capture}

The system should allow users to save information from multiple sources.

\textbf{Supported sources:}
\begin{itemize}[nosep]
    \item Web articles
    \item PDFs
    \item YouTube videos
    \item Manual notes
    \item Documentation pages
\end{itemize}

\textbf{For every saved source, the platform stores:}
\begin{itemize}[nosep]
    \item Title
    \item Content
    \item Source URL
    \item Date added
    \item Source type
    \item User metadata
\end{itemize}

\subsection{Pillar 2: Knowledge Organization}

The system automatically structures saved content.

Each item should be categorized using:

\textbf{Domain}
\begin{itemize}[nosep]
    \item Programming
    \item AI
    \item Databases
    \item System Design
    \item Networking
\end{itemize}

\textbf{Topic}
\begin{itemize}[nosep]
    \item Redis
    \item PostgreSQL
    \item Docker
    \item RAG
\end{itemize}

\textbf{Difficulty}
\begin{itemize}[nosep]
    \item Beginner
    \item Intermediate
    \item Advanced
\end{itemize}

\textbf{Content Type}
\begin{itemize}[nosep]
    \item Article
    \item Video
    \item PDF
    \item Note
\end{itemize}

\textbf{Status}
\begin{itemize}[nosep]
    \item Unread
    \item Reading
    \item Learned
    \item Mastered
\end{itemize}

The user may manually edit classifications.

\subsection{Pillar 3: Knowledge Retrieval}

Users should be able to search their knowledge base naturally.

\textbf{Example queries:}
\begin{itemize}[nosep]
    \item What do I know about Redis?
    \item Show all notes about authentication.
    \item Find resources related to system design.
    \item Which articles discuss caching?
\end{itemize}

\textbf{The platform should return:}
\begin{itemize}[nosep]
    \item Relevant sources
    \item Extracted snippets
    \item Related concepts
    \item AI-generated summaries
\end{itemize}

\subsection{Pillar 4: AI Question Answering}

Users can ask questions against their knowledge base.

\textbf{Examples:}
\begin{itemize}[nosep]
    \item Explain Redis based on my saved resources.
    \item Summarize everything I know about databases.
    \item Compare PostgreSQL and MongoDB using my notes.
\end{itemize}

\textbf{Responses must:}
\begin{itemize}[nosep]
    \item Use only saved content.
    \item Cite sources.
    \item Link back to original documents.
\end{itemize}

\textbf{If information is unavailable}, the system should clearly indicate:

\begin{quote}
    \textit{``This information does not exist in your knowledge base.''}
\end{quote}

\subsection{Pillar 5: Knowledge Connections}

The system should automatically identify relationships between content.

\textbf{Example:}

Article A:
\begin{itemize}[nosep]
    \item Redis caching
\end{itemize}

Article B:
\begin{itemize}[nosep]
    \item System design interview notes
\end{itemize}

Detected connection:
\begin{itemize}[nosep]
    \item Scalability
    \item Performance optimization
\end{itemize}

\textbf{Users should be able to explore:}
\begin{itemize}[nosep]
    \item Related documents
    \item Related topics
    \item Related concepts
\end{itemize}

\subsection{Pillar 6: Revision \& Retention}

The system should actively help users retain knowledge.

\textbf{Each knowledge item maintains:}
\begin{itemize}[nosep]
    \item Date added
    \item Last revision date
    \item Revision count
    \item Confidence score
    \item Knowledge score
\end{itemize}

\subsubsection{Revision Recommendation Engine}

The system recommends revisions based on:

\begin{itemize}[nosep]
    \item Time since learning
    \item Time since last revision
    \item Confidence level
    \item Topic importance
    \item User activity
\end{itemize}

\textbf{Example:}

Today's Revision List:
\begin{enumerate}[nosep]
    \item Redis Fundamentals
    \item Bellman-Ford Algorithm
    \item HTTP Caching
\end{enumerate}

\subsubsection{Smart Revision Material}

Instead of reopening entire documents, the system generates:

\begin{itemize}[nosep]
    \item Key points
    \item Summaries
    \item Flashcards
    \item Quiz questions
    \item Important concepts
\end{itemize}

All generated from saved knowledge.

\subsection{Pillar 7: Learning Progress Tracking}

The platform tracks learning activity.

\textbf{Metrics:}
\begin{itemize}[nosep]
    \item Total knowledge items
    \item Topics learned
    \item Revision streak
    \item Most studied domains
    \item Weakest knowledge areas
    \item Knowledge growth over time
\end{itemize}

\section{User Dashboard}

The home dashboard should display:

\subsection{Knowledge Overview}
\begin{itemize}[nosep]
    \item Total saved resources
    \item Total topics
    \item Total revisions completed
\end{itemize}

\subsection{Today's Revision Queue}

Recommended items for revision.

\subsection{Recently Added Knowledge}

Recent articles, notes, and PDFs.

\subsection{Learning Insights}

Examples:
\begin{itemize}[nosep]
    \item Networking has not been revised in 30 days.
    \item Databases is your most studied topic.
    \item You are currently learning Backend Development.
\end{itemize}

\section{Search Experience}

The platform should support:

\subsection{Structured Search}

Examples:
\begin{itemize}[nosep]
    \item Topic = Databases
    \item Status = Learned
    \item Difficulty = Intermediate
\end{itemize}

\subsection{Semantic Search}

Examples:
\begin{itemize}[nosep]
    \item How can I improve backend performance?
    \item Explain caching strategies.
\end{itemize}

\section{Knowledge Graph View (Future)}

Visual representation of relationships between:

\begin{itemize}[nosep]
    \item Topics
    \item Concepts
    \item Documents
\end{itemize}

\textbf{Example:}

\begin{center}
Redis \\
$\downarrow$ \\
Caching \\
$\downarrow$ \\
Performance \\
$\downarrow$ \\
Scalability
\end{center}

Users can explore how concepts connect across their knowledge base.

\section{Success Metrics}

The product is successful if users can:

\begin{enumerate}[nosep]
    \item Find previously learned information within seconds.
    \item Reduce repeated relearning.
    \item Build a structured long-term knowledge repository.
    \item Maintain a consistent revision habit.
    \item Discover new connections between saved knowledge.
\end{enumerate}

\section{Product Vision Statement}

Personal Knowledge OS transforms information into retained knowledge by combining intelligent organization, semantic retrieval, AI-assisted understanding, and personalized revision guidance within a single learning platform.

\end{document}
